% SPDX-License-Identifier: CC-BY-4.0
\documentclass[10pt]{article}

\usepackage[T1]{fontenc}
\usepackage[utf8]{inputenc}
\usepackage{lmodern}
\usepackage[margin=1in]{geometry}
\usepackage{amsmath,amssymb,amsthm}
\usepackage{enumitem}
\usepackage{fancyvrb}
\usepackage{booktabs}
\usepackage{tabularx}
\usepackage{array}
\usepackage{microtype}
\usepackage{parskip}
\usepackage{needspace}
\usepackage[hidelinks]{hyperref}
\usepackage{xurl}
\urlstyle{same}
\setlength{\emergencystretch}{2em}

\newtheorem{theorem}{Theorem}
\newtheorem{lemma}{Lemma}
\newtheorem{definition}{Definition}
\newtheorem{proposition}{Proposition}
\newcommand{\code}[1]{\texttt{#1}}
\newcommand{\Adv}{\mathsf{Adv}}
\newcommand{\Enc}{\mathsf{Enc}}
\newcommand{\Sign}{\mathsf{Sign}}
\newcommand{\Verify}{\mathsf{Verify}}
\newcommand{\Register}{\mathsf{RegisterFresh}}
\newcommand{\Consume}{\mathsf{ConsumeIfUnused}}
\newcommand{\ABIA}{\mathsf{ABIA}}
\DefineVerbatimEnvironment{protocol}{Verbatim}{fontsize=\footnotesize,frame=single,framesep=4pt}

\hypersetup{
  pdftitle={Action-Bound Injective Authorization under Signer Harvesting},
  pdfauthor={Iman Schrock},
  pdfsubject={Cryptographic protocols},
  pdfkeywords={authorization, digital signatures, replay, signer harvesting, formal verification}
}

\title{Action-Bound Injective Authorization\\under Signer Harvesting}
\author{Iman Schrock\\
EMILIA Protocol, Inc.\\
\texttt{team@emiliaprotocol.ai}\\
ORCID: \href{https://orcid.org/0009-0004-0290-5433}{0009-0004-0290-5433}}
\date{Preprint v3. August 2026.\\
\href{https://doi.org/10.5281/zenodo.21968577}{doi:10.5281/zenodo.21968577}}

\begin{document}
\maketitle

\begin{abstract}
A valid signature proves that a key signed a message. It does not, by itself,
prove that an authorization collector bound the signature to the action later
admitted, that two collected signatures belong to one approval ceremony, or that
the same authorization was admitted at most once. These gaps matter when the
collector renders requests, solicits signatures, and retains the audit record,
and is itself adversarial. We define Action-Bound Injective Authorization (ABIA),
a protocol property for this signer-harvesting setting. An approver signs a
closed encoding of an exact action and an independently issued authorization
instance. An enforcement domain admits the authorization only through a
linearizable consume-if-unused transition. We give explicit transplantation,
duplicate-admission, and quorum-substitution games. We prove that a successful
transplantation or quorum substitution yields a signature forgery or hash
collision, while duplicate admission yields an atomicity failure or enforcement
bypass. The result deliberately separates cryptographic authenticity from the
state-service assumption and does not claim exactly-once physical execution. We
also characterize four load-bearing assumptions by concrete attacks, state the
exact trace-prefix authenticity retained across a later key reveal, and prove
that ordinary signatures alone cannot support offline anti-backdating. We show
how Ed25519 and ML-DSA context strings instantiate the construction. Four
Tamarin models check the corresponding symbolic properties against an attacker
that can solicit honest signatures on arbitrary contexts. Nine deliberately
weakened variants produce concrete counterexample traces without weakening the
signature primitive.
\end{abstract}

\section{Introduction}

Per-transaction signatures are an old and useful idea. Payment cards,
hardware wallets, transaction-signing authenticators, and authorization tokens
all bind some request data to a cryptographic authenticator. The security question
studied here begins one layer later: what may an enforcement point infer when the
service that assembled the request, displayed it for approval, collected the
signature, and retained the resulting record is adversarial?

Suppose an operator can ask an honest approver to sign arbitrarily many valid
requests. This is stronger than an attacker that merely steals a bearer token and
different from an attacker that extracts the approver's key. The operator may
harvest genuine signatures, replay them, splice signatures from different
approval ceremonies, alter unsigned labels, or present the same authorization to
concurrent executors. EUF-CMA security rules out a new signature on a message that
was never signed. It does not define which application fields must be in that
message, whether two signatures refer to the same ceremony, or whether a valid
authorization may be admitted twice.

We call the missing protocol property ABIA, short for \emph{Action-Bound
Injective Authorization}. ``Action-bound'' means that every accepted approval signs
a closed, injective encoding of the exact action, policy commitment, initiator,
audience, authorization instance, validity data, and approver slot. ``Injective''
means that, within one enforcement domain, one authorization instance can precede
at most one admitted provider attempt. The second property is stateful. A
signature cannot make a distributed consume operation linearizable.

The construction uses only standard signatures and hashes. It is not a new
signature primitive. Its contribution is a protocol experiment and a reduction
that keep three failure classes separate:

\begin{enumerate}[leftmargin=*]
\item \textbf{Transplantation.} An accepted action was not the context signed by
an uncompromised required key.
\item \textbf{Duplicate admission.} The same authorization instance precedes two
provider-entry events in one enforcement domain.
\item \textbf{Quorum substitution.} A 2-of-2 admission succeeds without two
distinct pinned approver keys signing the same base authorization, or with the
initiator occupying an approver slot.
\end{enumerate}

The main theorem bounds the first and third events by signature unforgeability and
hash collision resistance. The second event is impossible only under complete
mediation by a linearizable register-and-consume interface. This division is the
point of the result. Systems often describe a signed token as ``single-use'' even
though its single-use property exists only in an enforcement service. ABIA names
that service and carries its failure probability as a separate term.

The paper contributes:

\begin{itemize}[leftmargin=*]
\item a signer-harvesting experiment in which the adversarial operator can obtain
valid signatures on adaptively chosen authorization contexts;
\item a closed protocol interface for single-approver and 2-of-2 authorization;
\item a reduction that separates exact-context authenticity from single-use
admission and from complete provider mediation;
\item necessity results for injective encoding, consume-if-unused state,
initiator binding and exclusion, and fresh non-recreatable registration;
\item a precise separation between trace-prefix authenticity after a later key
reveal and the stronger offline anti-backdating property that ordinary
signatures alone cannot provide;
\item symbolic confirmation in Tamarin, including negative models that expose
same-authorization replay, stale registry views, downgrade, authority expansion,
profile substitution, and action-key substitution when their checks are removed;
and
\item an explicit boundary: ABIA proves neither human identity nor rendering
fidelity, policy correctness, log completeness, or downstream physical effects.
\end{itemize}

\section{Attacks and Necessary Conditions}

The following attacks explain why ordinary signature security does not imply
ABIA. Every accepted signature in these traces is genuine and verifies. The
attacker changes only application encoding, ceremony binding, or mutable state.
Four traces are central to the proof, and a fifth shows how the same failure
appears in a generic approval workflow.

\subsection{Context transplant under non-injective encoding}

Suppose distinct typed contexts $x\neq x'$ have the same accepted byte encoding.
The operator asks an honest approver to sign $x$ and later presents the returned
signature with $x'$. The verifier reconstructs the same bytes, so the signature
verifies even though a security-relevant field changed.

\begin{protocol}
Operator -> Approver:  display x; request Sign(Enc(x))
Approver -> Operator:  sigma = Sign(sk, Enc(x))
Operator -> Executor:  present (x', sigma), x' != x
Executor:               Enc(x') = Enc(x); Verify = true; accept x'
\end{protocol}

This is not an EUF-CMA forgery because the accepted message bytes were queried to
the signing oracle. It is an application-language failure. A closed, injective
encoding prevents the equality above; collision resistance additionally prevents
two distinct canonical actions from sharing the action digest embedded in the
slot message.

\subsection{Replay without admission state}

An approver signs message $m$ once and the operator sends the same valid artifact
to two executors. Both verify it and enter the provider.

\begin{protocol}
Operator -> Approver:  request Sign(m)
Approver -> Operator:  (m, sigma)
Operator -> Executor1: (m, sigma)        Executor1 -> Provider: entry(u)
Operator -> Executor2: (m, sigma)        Executor2 -> Provider: entry(u)
\end{protocol}

Neither acceptance is a forgery. A nonce inside $m$ is insufficient unless shared
state remembers and atomically consumes it. Per-process caches fail when executors
race or restart. The 11-step Tamarin counterexample reproduces this trace by
removing only the linear \code{AdmissionAvailable} input.

\subsection{Quorum splicing without a shared instance}

Suppose two approver slots sign the same action and policy but no common
authorization instance. The operator starts ceremonies $C_1$ and $C_2$, obtains a
slot-1 signature from $C_1$ and a slot-2 signature from $C_2$, and combines them.

\begin{protocol}
Operator -> Approver1: Sign(action, policy, slot1, C1)
Approver1 -> Operator: sigma1
Operator -> Approver2: Sign(action, policy, slot2, C2)
Approver2 -> Operator: sigma2
Operator -> Executor:  (sigma1, sigma2) as one alleged quorum
\end{protocol}

Both signatures are genuine, but no two approvers joined one ceremony. A fresh
shared instance $q$, independently issued and included in both slot messages,
turns ceremony membership into byte equality checked by the verifier. Per-slot
nonces remain distinct and do not replace $q$.

\subsection{Initiator substitution without signed separation of duties}

Suppose self-approval is checked against an unsigned initiator label. The operator
collects genuine approvals under label $i$, changes the label to $i'$, and asks the
executor to apply separation of duties to the changed value.

\begin{protocol}
Operator -> Approvers: Sign(action, q, slots); unsigned initiator = i
Approvers -> Operator: sigma1, sigma2
Operator:              replace unsigned i by i'
Operator -> Executor:  (i', sigma1, sigma2); one approver = i'
Executor:               signatures verify under the substituted label
\end{protocol}

The historical Tamarin model found this attack directly. Binding $i$ into every
slot message makes relabeling a signed-message change. The executor must also
exclude the signed initiator from every approver slot; binding without exclusion
would faithfully authorize self-approval.

\subsection{Stale registry-view laundering}

An operator obtains valid approvals under an old authority view $h_0$. The
registry later publishes $h_1$, which revokes or narrows that authority. If the
executor checks only an operator-presented or unpinned view, the old evidence can
be evaluated as though $h_0$ were still current.

\Needspace{7\baselineskip}
\begin{protocol}
Registry:               publish h0; later publish h1 that revokes authority
Operator -> Approvers:  collect signatures bound to h0
Operator -> Executor:   present signatures plus stale view h0
Executor:               accept because no exact current-head check is required
\end{protocol}

The composed Tamarin comparison falsifies current-view acceptance in 20 steps.
The repair is not to distrust the signatures. It is to bind the relied-upon
authority material and require the executor's independently pinned current view.

\subsection{Generic approval-workflow transplant}

A common workflow records that a user clicked ``approve'' and then signs or logs
a token containing a request identifier, user identifier, and boolean decision.
The consequential action bytes remain in an operator-controlled database row.

\begin{protocol}
Operator -> User:       display action a and an Approve button
User -> Backend:        approve(request_id = R)
Backend -> Operator:    token = Sign(sk, R, user, approved)
Operator:               change database row R from a to a'
Operator -> Executor:   (R, a', token); signature verifies
\end{protocol}

This pattern loses the transplantation game immediately: the signed token never
commits to the canonical action the executor will submit. No vendor assumption is
needed. The minimal repair is to derive the action at the enforcement boundary,
place its canonical commitment and the policy, initiator, audience, and fresh
authorization instance inside the signed message, and atomically consume the
registered admission key.

\subsection{Nine machine-found separations}

The current artifacts retain nine deliberately weakened comparisons. Each changes
one relied-upon check while leaving the ideal signature algebra unchanged.

\small
\begin{tabularx}{\textwidth}{@{}>{\raggedright\arraybackslash}p{0.27\textwidth}>{\raggedright\arraybackslash}p{0.28\textwidth}>{\raggedright\arraybackslash}X@{}}
\toprule
Attack name & Omitted check & Machine-found consequence \\
\midrule
Focused Double Acceptance & checked acceptance consumes state & one genuine receipt is accepted twice in the focused model \\
Composed Double Entry & provider entry consumes shared state & one genuine authorization reaches provider entry twice after full composition \\
Stale Head Laundering & exact current registry head & revoked or narrowed authority is accepted under an old view \\
Class Downgrade & minimum assurance class & weaker evidence is accepted for a stronger class \\
Denial Becomes Authority & denial non-authorization & a signed denial is treated as positive authority \\
Scope Expansion & authority-scope pin & an action outside the signed scope is admitted \\
Profile Swap & complete profile pin & a different verification profile is substituted \\
Phantom Challenge & registration before exposure & an unregistered challenge is accepted \\
Action-Key Swap & canonical action-key binding & evidence for one action keys provider entry for another \\
\bottomrule
\end{tabularx}
\normalsize

The 11-step focused replay occurs in Single Approval. The 31-step no-consumption
trace and 20-step stale-view trace occur in Composed Reliance. The remaining six
paired traces occur in Enforcement Obligations in 14--16 steps. The two replay
lemmas expose the same missing state transition at different composition depths,
so the artifacts contain nine falsified comparisons but eight attack classes.
The initiator-relabeling attack is retained as repaired design history because
the strict current quorum model now binds the initiator.

\subsection{Necessity results}

The same traces give tightness results for the protocol assumptions.

\begin{proposition}[Injective encoding is necessary]
If the accepted authorization language admits distinct security contexts
$x\neq x'$ with identical signed bytes, then an operator with one honest
signature on $x$ wins transplantation for $x'$ with probability one.
\end{proposition}

\begin{proof}
Query the honest signer on the common byte string and present the returned
signature with $x'$. Verification succeeds on the same bytes. The contexts differ,
so transplantation holds, but EUF-CMA is not violated because that byte string was
queried.
\end{proof}

\begin{proposition}[Consume-if-unused is necessary]
If two provider-entry paths can verify the same authorization without sharing one
linearizable consume-if-unused transition, an operator can win duplicate admission
using one genuine receipt and no cryptographic break.
\end{proposition}

\begin{proof}
Deliver the receipt concurrently to the two paths. Both cryptographic checks are
identical and succeed. Without shared atomic state, neither path observes that the
other admitted the same key.
\end{proof}

\begin{proposition}[Signed initiator exclusion is necessary]
If the initiator is not bound into every slot message, an operator can relabel a
valid quorum before the separation-of-duties check. If the initiator is bound but
not excluded from approver slots, the protocol accepts self-approval directly.
\end{proposition}

\begin{proof}
The first case is the relabeling trace above. In the second, ask the signed
initiator to occupy a required slot and present the resulting valid signature;
all checks other than exclusion succeed.
\end{proof}

\begin{proposition}[Fresh non-recreatable registration is necessary]
If an operator may register an attacker-chosen stale instance, or recreate
$\textsc{unused}$ for a consumed admission key, then it can reuse previously
collected signatures to obtain a new successful consume and provider entry.
\end{proposition}

\begin{proof}
Keep the original signed slots. Reintroduce their shared instance or reset its
admission key to $\textsc{unused}$, then present the unchanged signatures. The
verifier reconstructs the original signed messages and the recreated state makes
the consume succeed. No forgery or collision is required.
\end{proof}

These attacks are minimal. They require no key extraction or failure of the
signature primitive. They differ only in protocol bindings and shared state,
which is why those assumptions appear as first-class objects in the definition
and theorem.

\section{Preliminaries and Assumptions}

\subsection{Signatures and encoding}

Let $\mathsf{Sig}=(\mathsf{Kg},\Sign,\Verify)$ be a digital signature scheme.
For each approver $j$, $(sk_j,pk_j)\leftarrow\mathsf{Kg}(1^\lambda)$. We use the
standard EUF-CMA experiment and write
$\Adv^{\mathrm{euf\mbox{-}cma}}_{\mathsf{Sig}}(\mathcal B)$ for the probability
that adversary $\mathcal B$ produces a valid signature on a message not submitted
to its signing oracle.

Let $H:\{0,1\}^*\rightarrow\{0,1\}^\kappa$ be collision resistant. Let
$\Enc$ be a deterministic, injective encoding on a closed typed schema. ``Closed''
means that a parser rejects unknown members, missing required members, duplicate
members, values outside the declared types, and non-canonical representations.
The proof needs injectivity, not a particular wire format. A conforming
implementation may use a canonical encoding such as deterministic CBOR or JCS
over a closed JSON profile, provided it establishes the same property.

Domain-separation labels are literal protocol constants and are included in every
encoded tuple. All equality checks below are byte equality after strict parsing;
display strings and workflow identifiers are not security inputs.

\subsection{Authorization context}

A canonical action $a$ contains an action type and every material parameter the
executor will use. Let $p$ be a commitment to the exact policy version, $i$ the
initiator identifier, $r$ the relying-party audience, and $q$ a fresh shared
authorization instance. The base context is

\[
  b=(a,p,i,r,q).
\]

For approver slot $v$, let $n_v$ be a fresh slot nonce and $w_v$ a closed validity
object containing issuance and expiration bounds. The slot context is
$x_v=(b,v,n_v,w_v)$ and the signed message is

\[
m_{b,v}=\Enc(\code{ABIA-SIGN-v1},H(\Enc(a)),p,i,r,q,v,n_v,w_v).
\]
We write $m(x_v)=m_{b,v}$ for this deterministic mapping from a well-typed
semantic slot context to its signed byte string.

The shared $q$ binds all slot signatures to one ceremony. Distinct $n_v$ values
keep the slot contexts distinct. The relying party or another trusted component
for enforcement domain $d$ issues $q$ and durably registers the corresponding
admission key

\[
u_b=H(\Enc(\code{ABIA-ADMIT-v1},d,H(\Enc(a)),p,i,r,q)).
\]

The adversarial operator may relay $q$. It cannot cause acceptance under an
unregistered $q$, choose a replacement $q$, or recreate a consumed registration.

\subsection{Linearizable admission service}

For each domain $d$, the admission service exports:

\begin{itemize}[leftmargin=*]
\item $\Register_d(u)$, which atomically creates state
$u\mapsto\textsc{unused}$ if no state for $u$ has ever existed; and
\item $\Consume_d(u)$, which atomically changes
$u\mapsto\textsc{unused}$ to $u\mapsto\textsc{admitted}$ and returns success, or
returns failure without entering a provider for every other state.
\end{itemize}

Operations are linearizable and durable. No transition recreates
$\textsc{unused}$ for the same key. \emph{Complete mediation} means that every
entry into a provider in $d$ is preceded by a successful $\Consume_d(u)$ for
the same request. These are system assumptions, not consequences of signature
security.

\section{Protocol}

The following algorithms define the analyzed 2-of-2 profile. The single-approver
profile uses one slot and the same admission interface.

\begin{protocol}
Issue_d(a, p, i, r):
    repeat q <- {0,1}^lambda
    u <- H(Enc("ABIA-ADMIT-v1", d, H(Enc(a)), p, i, r, q))
    until RegisterFresh_d(u) returns success
    return (b = (a,p,i,r,q), u)

Approve(sk_j, b, v_j, n_j, w_j):
    m_j <- Enc("ABIA-SIGN-v1", H(Enc(a)), p, i, r,
               q, v_j, n_j, w_j)
    return (v_j, n_j, w_j, Sign(sk_j, m_j))

VerifySlot(pk_j, b, v_j, n_j, w_j, sigma_j):
    reject unless all inputs parse under the closed schema
    reject unless w_j is current under the relying party's clock policy
    m_j <- Enc("ABIA-SIGN-v1", H(Enc(a)), p, i, r,
               q, v_j, n_j, w_j)
    return Verify(pk_j, m_j, sigma_j)

VerifyAndConsume_d(b, slots[1], slots[2]):
    reject unless q equals the independently issued instance
    reject unless both slots contain the same base context b
    reject unless ids, pinned keys, slots, and nonces are pairwise distinct
    reject if either approver id equals initiator i
    reject unless VerifySlot succeeds for both pinned keys
    u <- H(Enc("ABIA-ADMIT-v1", d, H(Enc(a)), p, i, r, q))
    reject unless ConsumeIfUnused_d(u) returns success
    emit ProviderEntry(d, u, a)
\end{protocol}

The executor constructs $a$ from the operation it is about to submit. It does not
accept an operator-supplied label as the action. The policy commitment $p$ is an
exact digest or equivalent immutable identifier, not merely a policy name. The
audience $r$ prevents a signature collected for one relying party from being
accepted by another relying party that uses different trust and policy inputs.

The protocol emits \code{ProviderEntry}, not \code{EffectOccurred}. If the
provider accepts a request and its acknowledgement is lost, the outcome is
indeterminate. The executor must reconcile using provider state or an idempotency
interface. ABIA does not authorize a blind retry.

\section{Signature Instantiations and Context Strings}

ABIA is signature-scheme agnostic, but two current instantiations clarify the
role of cryptographic context strings.

With Ed25519, the approver signs the complete $m_{b,v}$ byte string using the
Ed25519 algorithm in RFC~8032~\cite{rfc8032}. The literal
\code{ABIA-SIGN-v1} remains inside the injective application encoding. RFC~8032
also defines Ed25519ctx, but deployments need not depend on that variant because
many interfaces expose only pure Ed25519. In either case, the verifier reconstructs
the complete application message rather than signing an operator-provided digest
with an unsigned description beside it.

FIPS~204 standardizes ML-DSA and HashML-DSA and permits a context string of at
most 255 bytes~\cite{fips204}. An ML-DSA profile may set that string to the ASCII
bytes \code{ABIA-SIGN-v1} and require the verifier to use the same value. This
provides cryptographic domain separation. It does not, by itself, bind the
canonical action, policy, initiator, audience, ceremony instance, slot, nonce, or
validity object, and it carries no single-use state. Those fields still belong in
$m_{b,v}$, and \code{ConsumeIfUnused} remains necessary.

Thus the standardized context parameter and ABIA solve different problems. The
former separates uses of one signature primitive across protocols. The latter
defines the full typed statement authorized in this protocol and composes that
statement with admission state. Treating an ML-DSA context string as a substitute
for the application message would lose the transplantation game while every
signature continued to verify.

\section{Security Experiments}
\label{sec:experiments}

The challenger generates and pins all approver keys, initializes the admission
service for domain $d$, and gives every public key to adversary $\mathcal A$. The
adversary controls the network, operator, request transport, signature collector,
and artifact presenter. It may:

\begin{itemize}[leftmargin=*]
\item adaptively query $\mathsf{SignCtx}_j(x)$ for any well-typed slot context
$x$ under any uncompromised approver key $j$;
\item request issuance of honest authorization instances for chosen
$(a,p,i,r)$ tuples;
\item replay, suppress, reorder, and splice every returned artifact;
\item invoke verification and admission concurrently; and
\item reveal designated approver keys, after which no authenticity claim is made
for that key.
\end{itemize}

On a query $\mathsf{SignCtx}_j(x)$, the challenger strictly parses $x$, computes
the protocol message $m(x)$, records $\mathsf{Signed}(j,x,m(x))$, and returns
$\Sign(sk_j,m(x))$. Malformed or out-of-schema contexts are rejected. This
semantic-context oracle models signer harvesting without assuming that an honest
approver signs only a desirable action. The theorem asks a narrower question: if
a context is accepted under an unrevealed key, did that exact semantic context
appear in a prior signing event? Recording $x$ as well as $m(x)$ is essential:
otherwise a non-injective encoding transplant would be mislabeled as an ordinary
signing query on the same byte string.

\subsection{Trace events and corruption timing}

The experiment records the following ordered events:
\[
\begin{aligned}
&\mathsf{Issued}(d,b,u_b), &&
\mathsf{Signed}(j,x_v,m_{b,v}),\\
&\mathsf{Consumed}(d,u_b), &&
\mathsf{ProviderEntry}(d,u_b,a).
\end{aligned}
\]
An issued event occurs only after successful fresh registration of $u_b$ and
before any honest signing query for that authorization instance is answered. A
consumed event is recorded at the linearization point of a successful consume
operation. A provider-entry event is recorded immediately before crossing the
provider boundary and only after all cryptographic, policy, and state checks have
succeeded.

Key $j$ is unrevealed for an acceptance if no $\mathsf{Reveal}(j)$ event precedes
that acceptance. A later reveal does not retroactively erase the trace relation;
the theorem is about acceptance-time authenticity. Conversely, once a reveal has
occurred, the experiment makes no claim for later signatures under that key. This
is the conventional weak-corruption boundary needed by the per-key EUF-CMA
reduction. Forward-secure signatures could strengthen the post-compromise story,
but they are not assumed here.

The operator controls scheduling. It may interleave issuance, signing queries,
verification, and consume calls across arbitrarily many authorizations and may
delay a genuine signature until its validity window is nearly exhausted. The
challenger gives it no trusted session channel between approvers. Shared ceremony
membership therefore comes only from equality of the independently issued $q$
inside each signed context, not from network ordering or an operator-provided
session label.

The model treats a provider boundary as an abstract event. Complete mediation
requires every such event to be preceded by a successful consume on the same
$(d,u_b,a)$ relation. It does not require the provider to complete, acknowledge,
or produce the intended downstream effect. Those later events are outside the
ABIA experiment.

\subsection{Trace-prefix authenticity and the anti-backdating boundary}

The following corollary fixes the temporal scope of the authenticity experiment.
It is a trace-prefix statement, not a forward-secure logging or offline
anti-backdating property: a key reveal after an already recorded acceptance does
not change which exact context had to be signed before that acceptance.

\begin{lemma}[Trace-prefix authenticity]
Suppose an acceptance under pinned key $pk_j$ occurs at trace position $t_a$, no
reveal of $sk_j$ precedes $t_a$, and a reveal occurs later at $t_r>t_a$. Except
with probability bounded by
$\Adv^{\mathrm{euf\mbox{-}cma}}_{\mathsf{Sig}}(\mathcal B_j)+
\Adv^{\mathrm{cr}}_H(\mathcal C_H)$, a
$\mathsf{Signed}(j,x^*,m(x^*))$ event for the exact accepted slot context occurred
before $t_a$.
\end{lemma}

\begin{proof}
Stop the experiment at $t_a$. The later reveal is outside this prefix and cannot
affect its distribution. If the accepted message was not produced by a prior
semantic-context query before $t_a$, the pair is an EUF-CMA forgery for $j$. If
the same bytes were produced for another recorded context, injective outer
encoding and collision resistance make that context equal to the accepted
context. The later
reveal changes what the adversary can sign after $t_r$, but it cannot insert a
new signing event into the already ordered prefix.
\end{proof}

The Single Approval Tamarin model checks the symbolic counterpart in the named
later-reveal lemma: an acceptance before a later \code{RevealLtk} has prior
user-verification and signing events over the exact accepted context.

\begin{proposition}[Offline anti-backdating is impossible here]
Ordinary signatures and a signer-asserted timestamp do not let an offline verifier
distinguish a receipt genuinely signed before key reveal from one minted after
reveal with an earlier timestamp inside its signed message.
\end{proposition}

\begin{proof}
After learning $sk_j$, the adversary runs the honest signing algorithm on a
message containing any claimed earlier time. The resulting pair has the same
verification relation as a pair produced before compromise. An offline verifier
given only $(m,\sigma,pk_j)$ has no independent event ordering with which to
separate the two cases.
\end{proof}

A stronger anti-backdating claim needs an independently witnessed acceptance
before the reveal, such as inclusion in an append-only signed checkpoint with
trusted ordering, or a forward-secure signature scheme whose erased epoch keys
cannot be reconstructed after compromise. Neither mechanism is assumed in ABIA,
so the paper makes the prefix claim and not the stronger offline claim.

\begin{definition}[Transplantation]
$\mathcal A$ wins $\mathsf{Transplant}$ if the verifier accepts
$x^*=(b^*,v^*,n^*,w^*)$ under an unrevealed pinned key $j$, but no event
$\mathsf{Signed}(j,x^*,m(x^*))$ precedes that acceptance. A signing event for a
different semantic context $x\neq x^*$ does not defeat the winning condition,
even if a defective encoding maps $x$ and $x^*$ to the same bytes.
\end{definition}

\begin{definition}[Duplicate admission]
$\mathcal A$ wins $\mathsf{Duplicate}$ if two distinct provider-entry events in
the same domain $d$ are both labeled by the same admission key $u_b$.
\end{definition}

\begin{definition}[Quorum substitution]
$\mathcal A$ wins $\mathsf{QuorumSub}$ if a 2-of-2 provider entry occurs while at
least one of the following is true: fewer than two distinct pinned approver
identifiers and keys were accepted; a required unrevealed key has no prior
signing event for its exact accepted slot context; the two slots do not share
byte-identical $b$; or the
signed initiator occupies an approver slot.
\end{definition}

The combined experiment $\ABIA$ returns one if any of the three events occurs.
Let $\mathsf{AtomicityFail}_d$ denote any violation of the register-and-consume
contract, and let $\mathsf{MediationBypass}_d$ denote a provider entry that did not
pass through that contract.

\section{Security Argument}

\begin{lemma}[Exact-context binding]
Conditioned on no collision in $H$ and injectivity of $\Enc$, equality of two
valid slot messages implies equality of their action, policy commitment,
initiator, audience, authorization instance, slot, slot nonce, and validity
object.
\end{lemma}

\begin{proof}
The outer encodings carry the same domain-separation label and parse under one
closed schema. Injectivity therefore gives equality of every encoded field. In
particular, the action digests are equal. If the canonical action encodings were
different, they would form a collision in $H$; conditioned on no such collision,
the actions are equal. All remaining fields occur directly in the injective outer
encoding and are therefore equal.
\end{proof}

\begin{lemma}[Transplantation reduction]
For every transplantation adversary $\mathcal A$, there are EUF-CMA forgers
$\mathcal B_j$ and a collision finder $\mathcal C_H$ such that
\[
\Pr[\mathsf{Transplant}_{\mathcal A}=1]
\leq \sum_j \Adv^{\mathrm{euf\mbox{-}cma}}_{\mathsf{Sig}}(\mathcal B_j)
+\Adv^{\mathrm{cr}}_H(\mathcal C_H).
\]
\end{lemma}

\begin{proof}
The reduction is made explicit in the game sequence below. In short, an accepted
signature under unrevealed key $pk_j$ either has no prior semantic signing event,
in which case its message is an EUF-CMA forgery for $j$, or has a recorded query
on the same bytes. In the latter case the exact-context-binding lemma makes the accepted
context equal to the signed context unless a hash collision occurred. A union
bound over the pinned required keys and the collision event gives the inequality.
\end{proof}

\begin{lemma}[Injective admission]
Conditioned on linearizability, durable state, and complete mediation in domain
$d$, $\Pr[\mathsf{Duplicate}_{\mathcal A}=1]=0$.
\end{lemma}

\begin{proof}
By complete mediation, both provider entries must follow successful
$\Consume_d(u_b)$ operations. Linearizability places those calls in a total order
consistent with real time. The first success changes the only state for $u_b$ from
$\textsc{unused}$ to $\textsc{admitted}$. No allowed transition recreates
$\textsc{unused}$, so the second call cannot also succeed. A duplicate entry thus
implies either a state-service failure or a path around the service.
\end{proof}

\subsection{Game sequence}

We now give the hybrids underlying the concrete bound. Each game returns one if
the real ABIA experiment returns one, except where the game explicitly aborts.
The challenger records every semantic signing query as a triple $(j,x,m(x))$ and every issued
authorization as $(d,a,p,i,r,q,u_b)$. It also records the linearization point and
return value of every register and consume operation. This bookkeeping is not
visible to the adversary.

\paragraph{Game $G_0$: the real experiment.}
The challenger runs the protocol and adversary exactly as defined in
Section~\ref{sec:experiments}.
Thus
\[
  \Pr[G_0=1]=\Pr[\ABIA_{\mathcal A}=1].
\]

\paragraph{Game $G_1$: exclude action-hash collisions.}
The challenger aborts if two distinct canonical action encodings appearing in
recorded signing contexts or accepted contexts have the same hash. An algorithm $\mathcal C_H$
that runs $G_0$ and outputs the first such pair is a collision finder. Therefore
\[
  \left|\Pr[G_0=1]-\Pr[G_1=1]\right|
  \leq \Adv^{\mathrm{cr}}_H(\mathcal C_H).
\]
Conditioned on no abort, equality of action digests implies equality of canonical
action encodings. Together with injectivity of $\Enc$, equality of signed-message
bytes implies equality of every field in the slot context.

\paragraph{Game $G_2$: exclude accepted semantically unsigned contexts.}
For each unrevealed pinned key $j$, the challenger aborts if verification accepts
$(x^*,m^*,\sigma)$ under $pk_j$ while $(j,x^*,m^*)$ is absent from the
signing-query table. Under the $G_1$ condition and injectivity of the outer
encoding, absence of that triple also means that no query under $j$ produced
$m^*$: a different recorded context with the same message would contradict the
exact-context-binding lemma.
The reduction for one slot is the following explicit algorithm.

\begin{protocol}
Forger B_j(pk*, Sign*):
    set pk_j <- pk*; generate (sk_k, pk_k) for every k != j
    initialize the ABIA experiment and run A on all public inputs
    on A.SignCtx(j, x):
        reject unless x is well-typed; m <- m(x)
        record (j,x,m); return Sign*(m)
    on A.SignCtx(k, x), k!=j:
        reject unless x is well-typed; m <- m(x)
        record (k,x,m); return Sign(sk_k,m)
    on A.Reveal(k), k!=j:  return sk_k
    on A.Reveal(j):        abort this fixed-j reduction
    on accepted (x*,m*,sigma) under pk_j before Reveal(j):
        if (j,x*,m*) is absent from the query table:
            assert no row (j,x,m*) exists, by G1 and exact-context binding
            output (m*,sigma)
        otherwise continue the simulation
\end{protocol}

Registration, consumption, scheduling, and all public computations are simulated
exactly. The fixed-$j$ simulation is perfect up to either the first reveal of
$j$ or the accepted-unqueried event. A reveal before that event removes $j$ from
the experiment's authenticity claim. If the event occurs first, $\mathcal B_j$
stops immediately, so it never needs to reveal the EUF-CMA challenge secret. The
message was not queried to its EUF-CMA oracle and the signature verifies, which
is a valid forgery. This is exactly where the semantic query record matters: an
equal-byte query from a different context is ruled out by the encoding and hash
conditions rather than silently counted as authorization of $x^*$.

The standard multi-user union bound gives
\[
  \left|\Pr[G_1=1]-\Pr[G_2=1]\right|
  \leq \sum_j
  \Adv^{\mathrm{euf\mbox{-}cma}}_{\mathsf{Sig}}(\mathcal B_j).
\]
The sum includes only unrevealed keys required by the accepted policy. A key
reveal removes that key from the authenticity claim instead of asking the
reduction to simulate a forbidden reveal of its challenge secret.

\paragraph{Game $G_3$: replace accepted slots by recorded exact contexts.}
In $G_2$, every accepted slot under an unrevealed key has a matching recorded
triple $(j,x^*,m(x^*))$. The challenger now associates that slot with the
corresponding semantic query. This is a conceptual rewrite and changes no
probability. By the closed-schema checks, injectivity of $\Enc$, and the
no-collision condition from $G_1$, the recorded query and accepted slot agree on
the action, policy commitment, initiator, audience, authorization instance, slot,
slot nonce, and validity object. Transplantation is therefore impossible in
$G_3$.

\paragraph{Game $G_4$: condition on the admission contract.}
The challenger aborts on $\mathsf{AtomicityFail}_d$ or
$\mathsf{MediationBypass}_d$. This changes the success probability by at most
\[
  \Pr[\mathsf{AtomicityFail}_d\lor\mathsf{MediationBypass}_d].
\]
In $G_4$, every provider entry follows a successful consume operation. If two
entries carried the same $u_b$, their successful consumes would be ordered by
linearizability. The first moves the unique durable state for $u_b$ from
$\textsc{unused}$ to $\textsc{admitted}$; the second cannot succeed. Duplicate
admission is therefore impossible.

The verifier's deterministic distinctness, shared-base-context, pinned-key, and
initiator-exclusion checks remain. In $G_3$ every accepted unrevealed slot has an
exact recorded semantic signing event, so those checks also make quorum substitution
impossible. Hence $\Pr[G_4=1]=0$.

\begin{table}[ht]
\centering
\caption{Advantage accounting for the game sequence.}
\small
\begin{tabularx}{\textwidth}{@{}>{\raggedright\arraybackslash}p{0.15\textwidth}>{\raggedright\arraybackslash}p{0.28\textwidth}>{\raggedright\arraybackslash}p{0.24\textwidth}>{\raggedright\arraybackslash}X@{}}
\toprule
Transition & Bad event & Reduction or assumption & Probability charge \\
\midrule
$G_0\to G_1$ & distinct canonical actions share one digest & collision finder $\mathcal C_H$ & $\Adv_H^{\mathrm{cr}}(\mathcal C_H)$ \\
$G_1\to G_2$ & an unrevealed required key accepts a context without its exact signing event & one explicit forger $\mathcal B_j$ per required key & $\sum_j\Adv_{\mathsf{Sig}}^{\mathrm{euf\mbox{-}cma}}(\mathcal B_j)$ \\
$G_2\to G_3$ & associate accepted slots with recorded semantic contexts & injectivity and the $G_1$ condition & $0$ \\
$G_3\to G_4$ & atomic state fails or provider entry bypasses mediation & deployment contract & $\Pr[\mathsf{AtomicityFail}_d\lor\mathsf{MediationBypass}_d]$ \\
$G_4$ & any ABIA winning condition remains & deterministic verifier checks & $0$ \\
\bottomrule
\end{tabularx}
\normalsize
\end{table}

Each $\mathcal B_j$ makes exactly the signing-oracle queries that $\mathcal A$
makes for key $j$ and adds only key generation for the other slots, protocol
simulation, and table lookups. There is no guessing loss beyond the displayed
multi-user sum because the analysis defines one reduction per required key and
applies a union bound.

\subsection{Simulation and state boundaries}

The signature reduction does not simulate the admission service
cryptographically. Register and consume are ordinary oracle calls in the
experiment, and their failure probability remains an explicit term. This is
intentional. Replacing that term with a negligible cryptographic advantage would
silently assume the conclusion that concurrent single-use admission is correct.

The reduction also does not treat expiration, policy selection, identity proofing,
or display fidelity as signature properties. The verifier checks the closed
validity object and pinned policy inputs before acceptance, but the correctness of
those inputs belongs to the relying party. Likewise, a signature over action $a$
does not prove that a human was shown a faithful rendering of $a$. The experiment
only establishes which bytes were signed and which action the executor admitted.

\begin{lemma}[Quorum binding]
Conditioned on no signature forgery or hash collision, the 2-of-2 verification
algorithm cannot satisfy $\mathsf{QuorumSub}$ for unrevealed required keys.
\end{lemma}

\begin{proof}
The algorithm checks identifier, key, slot, and nonce distinctness and rejects the
signed initiator in either approver slot. It checks byte equality of the base
context before verifying one signature for each pinned slot key. If a required
slot context has no prior semantic signing event under that key, its accepted
signature is a forgery. If both events exist, exact-context binding makes both signed slot
contexts equal to the accepted contexts. The explicit base-context check then
makes both signatures authorize the same action, policy, initiator, audience, and
authorization instance. The explicit distinctness and initiator-exclusion checks
rule out the remaining winning conditions.
\end{proof}

\begin{theorem}[ABIA composition]
Let $\mathsf{Sig}$ be EUF-CMA secure, $H$ collision resistant, and $\Enc$
injective on the closed authorization schema. Consider any probabilistic
polynomial-time adversary $\mathcal A$ in the signer-harvesting experiment.
There are forgers $\mathcal B_j$ and a collision finder $\mathcal C_H$ such that
\[
\Pr[\ABIA_{\mathcal A}=1] \leq
\sum_j \Adv^{\mathrm{euf\mbox{-}cma}}_{\mathsf{Sig}}(\mathcal B_j)
+\Adv^{\mathrm{cr}}_H(\mathcal C_H)
+\Pr[\mathsf{AtomicityFail}_d\lor\mathsf{MediationBypass}_d].
\]
For the 2-of-2 profile, the sum ranges over the two unrevealed pinned approver
keys required by the accepted policy.
\end{theorem}

\begin{proof}
Games $G_0$ and $G_1$ differ only on an action-hash collision. Games $G_1$ and
$G_2$ differ only on an accepted context under an unrevealed pinned key without
an exact prior semantic signing event. The reductions above bound
those differences by the collision and per-key EUF-CMA advantages. Game $G_3$ is
identical to $G_2$ but uses exact-context binding to name the matching semantic
signing events, so transplantation and the signature-dependent quorum conditions cannot
occur. Game $G_4$ additionally conditions on successful atomic admission and
complete mediation, which excludes duplicate admission. The verifier's remaining
quorum conditions are deterministic checks. Thus $G_4$ has zero success
probability, and summing the game differences proves the inequality.
\end{proof}

\begin{theorem}[2-of-2 quorum corollary]
Suppose the 2-of-2 verifier admits base authorization $b$ with slot records
$x_1,x_2$ under distinct pinned keys $pk_1,pk_2$, neither key was revealed, and
neither approver identifier equals the signed initiator. Except with probability
bounded by
\[
  \Adv^{\mathrm{euf\mbox{-}cma}}_{\mathsf{Sig}}(\mathcal B_1)
  +\Adv^{\mathrm{euf\mbox{-}cma}}_{\mathsf{Sig}}(\mathcal B_2)
  +\Adv^{\mathrm{cr}}_H(\mathcal C_H),
\]
both pinned keys previously signed their assigned slot contexts over the same
action, policy commitment, initiator, audience, and authorization instance. No
single key or signed initiator fills both slots.
\end{theorem}

\begin{proof}
Apply Games $G_1$ through $G_3$ to the two required slots. In $G_2$, an accepted
slot without a matching semantic signing event is a forgery under that slot's
key. In $G_3$, each matching event binds the exact accepted slot fields. The verifier
checks byte equality of $b$, distinct slot identifiers, distinct approver
identifiers, distinct pinned keys, distinct nonces, and exclusion of the signed
initiator. Those deterministic checks give the stated 2-of-2 property. The
admission-service term is absent because this corollary concerns quorum
authenticity, not single-use provider entry.
\end{proof}

\subsection{Concrete loss and interpretation}

Let $N$ be the number of unrevealed pinned keys that may be required by an
accepted policy. If every per-key EUF-CMA advantage is at most
$\epsilon_{\mathrm{sig}}$, the collision advantage is at most
$\epsilon_H$, and the measured or assumed probability of either admission
failure event is at most $\epsilon_{\mathrm{state}}$, the theorem gives
\[
  \Pr[\ABIA_{\mathcal A}=1]
  \leq N\epsilon_{\mathrm{sig}}+\epsilon_H+\epsilon_{\mathrm{state}}.
\]
The signature reduction preserves the adversary's signing-query count for the
embedded key and adds only the cost of protocol simulation and table lookups.
There is no hidden claim that $\epsilon_{\mathrm{state}}$ is negligible. It is a
deployment term whose justification must come from the implementation of the
shared admission domain, including its concurrency, durability, and bypass
controls.

The bound is deliberately conservative in two ways. First, it uses a union bound
over required keys even when policy structure may make simultaneous failures
necessary. Second, it charges any detected action-hash collision once, regardless
of whether the collision would satisfy all remaining acceptance checks. Those
choices keep the reduction independent of one signature algorithm and one policy
cardinality. The price is a simple bound rather than a tight system-specific one.

The theorem gives at most one admitted provider attempt per authorization inside
$d$. It does not establish that the provider performed a physical effect, that an
effect occurred exactly once, or that a second independent domain will reject the
same authorization. Cross-domain single use requires a shared linearizable
admission service or an equivalent coordination protocol.

\section{Symbolic Analysis}

The computational proof isolates the assumptions. Four Tamarin models check the
same bindings under a Dolev-Yao network attacker. The models use Tamarin's
signing theory, so verification follows from the signing equations rather than
from an axiom saying that a signature is valid. The adversary controls the
network, learns all public keys, may reveal a device key through an explicit
rule, and may ask an honest approver to sign arbitrary attacker-chosen contexts.

User verification is a linear precondition to signing. Admission availability is
also a linear fact and checked acceptance consumes it. No trace restriction states
that acceptance is unique. This use of linear state follows the security-API
tradition: mutable authorization state is part of the model, not an informal note
outside it.

\subsection{Four models and checked obligations}

The filenames remain available with the artifact, but we use descriptive names
in the paper so that the result does not depend on project terminology.

\small
\begin{tabularx}{\textwidth}{@{}>{\raggedright\arraybackslash}p{0.18\textwidth}>{\raggedright\arraybackslash}X>{\raggedright\arraybackslash}p{0.18\textwidth}>{\raggedright\arraybackslash}p{0.20\textwidth}@{}}
\toprule
Model & Question & Strict result & Unsafe comparison \\
\midrule
Single Approval & Does acceptance imply a prior user-verified signature over the exact context, remain bound before a later reveal, and is checked acceptance injective? & 1 witness and 4 all-traces lemmas verified & replay without consumption, falsified in 11 steps \\
Quorum & Does 2-of-2 require two distinct pinned signers over the same action and ceremony, excluding the initiator? & 1 witness and 4 all-traces lemmas verified & historical initiator-relabeling trace when initiator binding was absent \\
Composed Reliance & Do exact action, profile, audience, authority view, issuer, quorum, and consumption remain bound through provider entry? & 1 witness and 9 all-traces lemmas verified & replay in 31 steps; stale registry view in 20 steps \\
Enforcement Obligations & Do downgrade refusal, denial non-authorization, scope and profile pinning, challenge consumption, and action-keyed entry hold together? & 2 witnesses and 8 all-traces lemmas verified & 6 weakened lemmas falsified in 14--16 steps \\
\bottomrule
\end{tabularx}
\normalsize

Across the four models, 25 all-traces obligations and five executable witnesses
verify. Nine deliberately weakened current-model comparisons produce attack
traces. The earlier quorum falsification is retained as design history but is not
counted among those nine current-model comparisons.

\subsection{Symbolic-to-computational correspondence}

The symbolic and computational analyses are not independent feature lists. The
table identifies which Tamarin state or event realizes each object in the games.

\small
\begin{tabularx}{\textwidth}{@{}>{\raggedright\arraybackslash}p{0.29\textwidth}>{\raggedright\arraybackslash}p{0.29\textwidth}>{\raggedright\arraybackslash}X@{}}
\toprule
Tamarin object & Computational object & Shared security meaning \\
\midrule
\code{Human\_Sign\_Receipt}, \code{Signed} & $\mathsf{SignCtx}_j(x)$ query and semantic query-table row & the attacker may harvest genuine signatures on chosen contexts \\
\code{Eq(verify(...),true)} under \code{PinnedKey} & $\Verify(pk_j,m,\sigma)$ under a pinned slot key & accepted bytes verify for the required key \\
\code{UV} before \code{Signed} & approver-side signing precondition & user verification gates signing, but does not prove rendering fidelity \\
\code{RevealLtk} ordered before or after \code{Accept} & weak-corruption boundary at acceptance & later reveal does not alter the accepted prefix; earlier reveal removes authenticity \\
\code{AdmissionAvailable} as a linear fact & successful $\Register_d(u)$ with state $\textsc{unused}$ & one fresh authorization instance has one consumable admission state \\
\code{AcceptChecked} consuming that fact & successful $\Consume_d(u)$ before provider entry & injectivity comes from mutable state, not signature algebra \\
Core authenticity and cross-action lemmas & Games $G_1$--$G_3$ & exact accepted context has a prior semantic signing event \\
Later-reveal prefix lemma & trace-prefix authenticity lemma & a later key reveal does not rewrite the already recorded trace prefix \\
Injective checked-acceptance lemma & Game $G_4$ & one registered key precedes at most one checked admission \\
Unchecked-injectivity comparison & $\mathsf{Duplicate}$ without the state contract & removing consumption yields a concrete replay \\
Quorum commit and distinct-slot lemmas & $\mathsf{QuorumSub}$ and its corollary & two distinct pinned keys sign one byte-identical base authorization, excluding the initiator \\
\bottomrule
\end{tabularx}
\normalsize

The correspondence has one deliberate asymmetry. Injectivity of the application
encoding and computational hardness of collisions are assumptions in the symbolic
model's term algebra, while Games $G_0$--$G_3$ expose and charge them explicitly.
Conversely, Tamarin explores arbitrary symbolic interleavings, while the
computational theorem carries state failure and mediation bypass as external
probability terms.

\subsection{Single approval and quorum}

In Single Approval, the request rule creates one
\code{AdmissionAvailable} fact for a fresh context. Checked acceptance consumes
it. The exact tool results are:

\begin{protocol}
executable_honest_receipt (exists-trace): verified (9 steps)
core_authenticity_uv_gated (all-traces): verified (11 steps)
acceptance_prefix_integrity_after_later_reveal (all-traces): verified (12 steps)
no_replay_across_actions (all-traces): verified (11 steps)
injective_acceptance_with_consumption (all-traces): verified (11 steps)
unchecked_acceptance_is_injective (all-traces): falsified - found trace (11 steps)
\end{protocol}

The unsafe rule changes only one thing: it accepts without consuming the linear
fact. The attacker delivers one genuine signature twice. No signature forgery or
key reveal occurs. This is the symbolic counterpart of the admission-service term
in Theorem 1.

Quorum enrolls two distinct approver identities under two pinned keys. Each
signature binds the action, initiator, shared authorization instance, slot, and
nonce. Commitment requires both signatures over the same base context and excludes
the initiator. Its five current lemmas verify in 4--27 steps:

\begin{protocol}
executable_quorum (exists-trace): verified (13 steps)
quorum_requires_two_distinct_uv_gated_signatures (all-traces): verified (27 steps)
initiator_cannot_self_approve (all-traces): verified (4 steps)
no_single_signer_fills_quorum (all-traces): verified (4 steps)
commit_requires_signature_over_that_action (all-traces): verified (8 steps)
\end{protocol}

An earlier revision
omitted the initiator from the signed tuple. Tamarin then found a trace in which
two honest signatures collected under one initiator label were committed under
another. Adding the initiator to the signed context restored the original lemmas;
no lemma was weakened.

\subsection{Composed models}

Composed Reliance places challenge issuance, exact action identity, profile and
audience, initiator, two distinct approvals, scoped authority, an exact registry
view, revocation status, issuer pinning, consumption, and provider entry in one
trace. Its strict path includes the following results:

\begin{protocol}
execution_requires_full_composition: verified (101 steps)
no_issuer_laundering: verified (785 steps)
no_cross_action_profile_or_audience_replay: verified (41 steps)
execution_has_honest_approvals_or_prior_compromise: verified (178 steps)
injective_execution_with_consumption: verified (250 steps)
\end{protocol}

Seven additional sanity or binding lemmas also verify. Removing consumption
produces a 31-step duplicate-entry trace. Removing exact registry-head binding
produces a 20-step stale-view trace.

Enforcement Obligations isolates six relied-upon checks and two sanity
properties. The strict model verifies downgrade refusal, signed-denial
non-authorization, authority-scope pinning, complete profile pinning,
registration-before-exposure and one-time challenge use, canonical action-keyed
provider entry, and fail-closed reservation. Six paired rules each omit one check.
Tamarin finds the corresponding downgrade, denial-as-authorization, authority
expansion, profile substitution, unregistered-challenge, and action-key
substitution traces in 14--16 steps.

These composed models are confirmation and separation tests, not additional
cryptographic theorems. They show that the declared bindings survive symbolic
composition and that removing each load-bearing check admits an attack while the
signature algebra stays ideal.

\subsection{Mechanized-analysis boundary}

The four models abstract WebAuthn internals, canonical parsers, natural-person
identity, trusted rendering, clocks, database durability, directory publication,
transparency-log completeness, and provider effects. A symbolic exact registry
view does not prove how that view was published. A linear fact models an atomic
state transition but does not prove a particular database implementation. The
models establish properties of the specified protocol under these declared
abstractions, not security of a deployment.

\section{Relation to Prior Work}

\paragraph{FIDO2 and WebAuthn security.}
Barbosa, Boldyreva, Chen, and Warinschi give a modular computational analysis of
WebAuthn and CTAP2 and prove authentication properties under their stated client,
token, and corruption models~\cite{fido2provable}. Guan et al. model WebAuthn and
CTAP2 together in ProVerif, while Schrempp gives a Tamarin model with an explicit
human role~\cite{fido2symbolic,fido2human}. Those works analyze authentication,
credential use, and the authenticator-client-relying-party ceremony. ABIA starts
after that interface: it assumes an enrolled signing key and asks what an executor
may infer when the operator that chooses, renders, collects, stores, and presents
action approvals is adversarial. Prior FIDO2 ceremony results do not define
single-use provider admission for an action authorization. The flat delta is:
ABIA makes the operator the signature-harvesting adversary and separates exact
signed-action authenticity from stateful single-use admission.

\paragraph{Transaction authorization.}
Work on smart-card payment protocols already distinguishes cryptographic proof of
critical transaction data from a reusable authentication event and observes that
publicly verifiable signatures provide different dispute evidence than shared-key
MACs~\cite{herreweghen99}. Hardware-wallet designs such as Notary focus on trusted
transaction display and correct device behavior~\cite{notary20}. The earlier
``see what you sign'' line makes trustworthy input and output part of binding
electronic signatures to human intent~\cite{weber98}. ABIA assumes no solution to
that display problem. Its narrower contribution is the collector-adversarial
experiment and the separation between exact signed context, ceremony binding, and
stateful single-use admission after signatures have been collected.

\paragraph{Security APIs and mutable protocol state.}
Formal analyses of PKCS\#11 show why mutable global state belongs inside a
security model: API safety can change when key attributes and permitted operations
change over time~\cite{pkcs11state}. ABIA uses the same broad methodological
lesson, although for a different goal. The linear admission fact is not a
cryptographic primitive. It represents the state transition that makes an
otherwise replayable valid authorization single-use. The positive and negative
Tamarin models keep the signature theory fixed and change only that state input.

\paragraph{Policy-aware signatures.}
Policy-Compliant Signatures make policy satisfaction part of an enhanced
signature primitive and provide privacy and composable-security notions for
attribute policies~\cite{pcs21}. Attribute-based signatures similarly hide the
signer's identity behind an attribute predicate~\cite{abs08}. ABIA provides
neither attribute privacy nor a new signature primitive. It uses ordinary
signatures under explicitly pinned keys and studies replay, ceremony splicing,
initiator substitution, and single-use admission by an external state service.
These are orthogonal properties. A policy-compliant or attribute-based signature
could replace the ordinary slot signature, but the admission theorem would still
require the state-service assumption.

\paragraph{Capabilities and least privilege.}
SPKI/SDSI, PolicyMaker, and KeyNote established public-key authorization and trust
management for distributed systems~\cite{spki,keynote}. Macaroons add contextual
caveats and decentralized attenuation to bearer credentials~\cite{macaroons}.
StatefulAuth allows client-defined authorization policies to depend on protected
request history~\cite{statefulauth}. These systems primarily decide what a
principal or bearer token may do. ABIA asks a different evidence question: whether
an accepted action is exactly the action signed by required approver keys under a
collector that may harvest honest signatures, and whether one such authorization
preceded at most one provider entry. StatefulAuth is especially close in
recognizing that authorization properties can require protected state. ABIA makes
the stronger atomicity and complete-mediation assumptions explicit because a
history integrity tag alone does not serialize concurrent consumption.

\paragraph{Formal protocol analysis.}
Tamarin and related symbolic tools have been used extensively for authentication
and key-exchange protocols~\cite{tamarin,dy}. Our use of Tamarin is conventional.
The contribution is the adversary query pattern and the paired positive and
negative models: the attacker may collect signatures on arbitrary contexts, and
the same ideal signing theory is used on both sides of each separation result.

\paragraph{Secure formats and semantic messages.}
Comparse treats concrete protocol formats as part of the security argument and
lifts bit-string assumptions to typed messages whose encodings satisfy explicit
format properties~\cite{comparse}. ABIA relies on the same core lesson: the
semantic authorization context must remain visible to the experiment, and a
byte-level signing query is not enough to define transplantation. We assume a
closed injective format rather than contributing a new verified parser framework.

\paragraph{Transparency, audit logs, and temporal evidence.}
Transparency overlays formalize append-only, externally checkable views using
signatures and dynamic list commitments, while forward-secure audit-log schemes
protect whole logs and verifiable excerpts after later compromise
~\cite{transparencyoverlays,secureauditlogs}. Those mechanisms can supply the
independent ordering evidence that ABIA intentionally omits. The trace-prefix
lemma only conditions on the experiment's already ordered events; it must not be
read as a transparency, secure-logging, or offline anti-backdating construction.

\paragraph{Web authorization and replay.}
The formal analysis of the OpenID FAPI 2.0 family shows how strong web attackers,
sender-constrained tokens, and replay-relevant protocol details belong in an
authorization proof tied to a standardization process~\cite{fapi2formal}. ABIA
studies a different boundary: an operator that can solicit genuine signatures on
many semantic action contexts. Its state term concerns atomic single-use provider
admission, not OAuth token presentation.

\section{Limitations}

ABIA proves that an unrevealed enrolled key signed an exact context. It does not
prove that the key belongs to a particular natural person. That claim depends on
directory governance and identity proofing. Pairwise-distinct enrolled identifiers
also do not prove two independent human wills.

The operator may render action $X$ while causing the signing device to sign action
$Y$. ABIA makes $Y$ recoverable and prevents the resulting signature from being
used for $Z$, but it cannot prove what the approver perceived. High-consequence
deployments need a trusted or independently authored display path. This is a
separate assumption, not a consequence of canonical encoding.

The theorem assumes the executor constructs the action from the operation it will
submit and that every provider entry passes through the admission service. An
operator able to bypass that executor is outside the enforcement claim and appears
as $\mathsf{MediationBypass}_d$ in the bound.

The proof covers one domain and the stated single-approver and 2-of-2 profiles. It
does not cover arbitrary threshold policies, key-directory publication,
revocation history, transparency-log consistency, clock synchronization, or
provider-side effects. A lost provider acknowledgement produces an indeterminate
outcome. Retrying safely requires provider reconciliation or a provider-enforced
idempotency key.

\section{Conclusion}

A signed authorization has two different security obligations. The signature and
closed encoding bind the authorization to an exact context. A linearizable,
completely mediated state transition binds that authorization to at most one
admitted provider attempt. Neither obligation implies the other. ABIA formalizes
their composition against an operator that may harvest genuine signatures on
arbitrary contexts. Its proof reduces context substitution to standard
cryptographic failures and leaves duplicate admission honestly attached to the
state service that enforces it. The Tamarin counterexamples show the same
boundary operationally: ideal signatures do not stop replay when consumption is
removed, and they do not stop initiator relabeling when the initiator is not
signed.

\section*{Artifact Availability}

This manuscript is archived as version v3 of the single-use authorization-receipt
preprint series at
\href{https://doi.org/10.5281/zenodo.21968577}{Zenodo record 21968577}.
It succeeds the shorter symbolic-analysis version at
\href{https://doi.org/10.5281/zenodo.21520973}{Zenodo record 21520973}; the
all-versions concept DOI is
\href{https://doi.org/10.5281/zenodo.21520972}{10.5281/zenodo.21520972}.

The four Tamarin theories, replay traces, exact tool summaries, and pinned rerun
scripts are available in the public repository:
\url{https://github.com/emiliaprotocol/emilia-protocol}. The theory files are:

\begin{protocol}
formal/tamarin/ep_receipt_core.spthy
formal/tamarin/ep_quorum_core.spthy
formal/tamarin/ep_reliance_composed.spthy
formal/tamarin/ep_six_claim_composed.spthy
\end{protocol}

The protocol definition, games, and proofs in this paper are self-contained; no
implementation or Internet-Draft is required to evaluate the result.

\begin{thebibliography}{99}

\bibitem{gmr88}
S. Goldwasser, S. Micali, and R. Rivest,
``A Digital Signature Scheme Secure Against Adaptive Chosen-Message Attacks,''
\emph{SIAM Journal on Computing}, 17(2), 1988.

\bibitem{dy}
D. Dolev and A. C. Yao,
``On the Security of Public Key Protocols,''
\emph{IEEE Transactions on Information Theory}, 29(2), 1983.

\bibitem{rfc8032}
S. Josefsson and I. Liusvaara,
``Edwards-Curve Digital Signature Algorithm (EdDSA),''
RFC 8032, January 2017.
\url{https://www.rfc-editor.org/rfc/rfc8032}.

\bibitem{fips204}
National Institute of Standards and Technology,
``Module-Lattice-Based Digital Signature Standard,''
FIPS 204, August 2024.
\url{https://doi.org/10.6028/NIST.FIPS.204}.

\bibitem{fido2provable}
M. Barbosa, A. Boldyreva, S. Chen, and B. Warinschi,
``Provable Security Analysis of FIDO2,''
\emph{CRYPTO 2021, Part III}, LNCS 12827, pp. 125--156, 2021.
\url{https://doi.org/10.1007/978-3-030-84252-9_5}.

\bibitem{fido2symbolic}
J. Guan, H. Li, H. Ye, and Z. Zhao,
``A Formal Analysis of the FIDO2 Protocols,''
\emph{ESORICS 2022}, LNCS 13556, pp. 3--21, 2022.
\url{https://doi.org/10.1007/978-3-031-17143-7_1}.

\bibitem{fido2human}
L. Schrempp,
``Formal Verification of FIDO2 with Human Interaction,''
Master's thesis, ETH Zurich, 2023.
\url{https://doi.org/10.3929/ethz-b-000610315}.

\bibitem{pkcs11state}
S. Delaune, S. Kremer, and G. Steel,
``Formal Security Analysis of PKCS\#11 and Proprietary Extensions,''
\emph{Journal of Computer Security}, 18(6), pp. 1211--1245, 2010.
\url{https://doi.org/10.3233/JCS-2009-0394}.

\bibitem{tamarin}
S. Meier, B. Schmidt, C. Cremers, and D. Basin,
``The TAMARIN Prover for the Symbolic Analysis of Security Protocols,''
\emph{CAV 2013}, LNCS 8044, pp. 696--701, 2013.

\bibitem{herreweghen99}
E. Van Herreweghen and N. Ferguson,
``Security Mechanisms in EMV,''
\emph{USENIX Workshop on Smartcard Technology}, 1999.

\bibitem{notary20}
A. Athalye, A. Belay, M. F. Kaashoek, R. Morris, and N. Zeldovich,
``Notary: A Device for Secure Transaction Approval,''
\emph{27th ACM Symposium on Operating Systems Principles}, pp. 97--113, 2019.
\url{https://doi.org/10.1145/3341301.3359661}.

\bibitem{weber98}
A. Weber,
``See What You Sign: Secure Implementations of Digital Signatures,''
\emph{5th International Conference on Intelligence and Services in Networks},
LNCS 1430, pp. 509--520, 1998.
\url{https://doi.org/10.1007/BFb0056995}.

\bibitem{pcs21}
C. Badertscher, C. Matt, and H. Waldner,
``Policy-Compliant Signatures,''
\emph{TCC 2021}, LNCS 13043, pp. 299--329, 2021.
\url{https://eprint.iacr.org/2021/1234}.

\bibitem{abs08}
H. Maji, M. Prabhakaran, and M. Rosulek,
``Attribute-Based Signatures: Achieving Attribute-Privacy and
Collusion-Resistance,'' Cryptology ePrint Archive, Paper 2008/328, 2008.

\bibitem{spki}
C. M. Ellison, B. Frantz, B. Lampson, R. Rivest, B. Thomas, and T. Ylonen,
``SPKI Certificate Theory,'' RFC 2693, 1999.

\bibitem{keynote}
M. Blaze, J. Feigenbaum, and A. D. Keromytis,
``KeyNote: Trust Management for Public-Key Infrastructures,''
RFC 2704, 1999.

\bibitem{macaroons}
A. Birgisson, J. G. Politz, U. Erlingsson, A. Taly, M. Vrable, and M. Lentczner,
``Macaroons: Cookies with Contextual Caveats for Decentralized Authorization in
the Cloud,'' \emph{NDSS}, 2014.

\bibitem{statefulauth}
L. Cao, L. Meng, D. Stefan, and E. Fernandes,
``Stateful Least Privilege Authorization for the Cloud,''
\emph{33rd USENIX Security Symposium}, pp. 3477--3495, 2024.

\bibitem{comparse}
T. Wallez, J. Protzenko, and K. Bhargavan,
``Comparse: Provably Secure Formats for Cryptographic Protocols,''
\emph{ACM CCS}, 2023.
\url{https://doi.org/10.1145/3576915.3623201}.

\bibitem{transparencyoverlays}
M. Chase and S. Meiklejohn,
``Transparency Overlays and Applications,''
\emph{ACM CCS}, 2016.
\url{https://eprint.iacr.org/2016/915}.

\bibitem{secureauditlogs}
G. Hartung,
``Secure Audit Logs with Verifiable Excerpts,''
\emph{CT-RSA}, 2016.
\url{https://doi.org/10.1007/978-3-319-29485-8_11}.

\bibitem{fapi2formal}
P. Hosseyni, R. K\"usters, and T. W\"urtele,
``Formal Security Analysis of the OpenID FAPI 2.0 Family of Protocols:
Accompanying a Standardization Process,''
\emph{ACM Transactions on Privacy and Security}, 2024.
\url{https://eprint.iacr.org/2024/1540}.

\end{thebibliography}

\end{document}
